% ==========================
% ПРЕАМБУЛА КОНСПЕКТА ЛЕКЦИЙ
% ==========================

% 1. НАЧАЛО: Языки и шрифты (СТРОГО В НАЧАЛЕ!)
\usepackage{cmap}            % Поиск и копирование русского текста из PDF
\usepackage[utf8]{inputenc}  % Кодировка исходного текста
\usepackage[T2A]{fontenc}    % Поддержка кириллических шрифтов
\usepackage[russian]{babel}  % Правила русского языка и переносы строк

% 2. ГЕОМЕТРИЯ: Размеры полей страницы
\usepackage[
    left=2cm,
    right=2cm,
    top=2cm,
    bottom=2cm
]{geometry}

% 3. МАТЕМАТИКА: Базовые пакеты и расширения
\usepackage{mathtools}       % Улучшения amsmath
\usepackage{amssymb}         % Классические мат. символы и шрифты
\usepackage{amsthm}          % Родной классический пакет теорем

% Знаки неравенств
\renewcommand{\le}{\leqslant}
\renewcommand{\leq}{\leqslant}
\renewcommand{\ge}{\geqslant}
\renewcommand{\geq}{\geqslant}

% Знак разности множеств
\renewcommand{\setminus}{\backslash}

% Операторы для отображений
\DeclareMathOperator{\Def}{Def}
\DeclareMathOperator{\D}{D}
\DeclareMathOperator{\Gr}{Gr}
\AtBeginDocument{
    \renewcommand{\Im}{\operatorname{Im}}
}

% 4. ГРАФИКА: Иллюстрации и рисунки
\usepackage{graphicx}        % Вставка готовых картинок (PNG, JPG, PDF)
\usepackage{tikz}            % Рисование графиков и схем кодом

% 5. ТИПОГРАФИКА: Оформление текста
\usepackage{indentfirst}     % Принудительный отступ (красная строка) для всех абзацев
\usepackage{microtype}       % Выравнивание краев текста без разрывов слов
\usepackage{enumitem}
\setlist{
    nosep,                   % Убирает вертикальные отступы между элементами списка
    topsep=3pt,              % Минимальный аккуратный отступ перед началом всего списка
    partopsep=0pt,
    parsep=0pt,
    leftmargin=*,            % Выравнивает маркеры по левому краю текста
    labelindent=\parindent   % Делает отступ списка равным обычной красной строке
}

% 6. ТЕОРЕМЫ И ОПРЕДЕЛЕНИЯ: Специальные блоки для математики
\usepackage{tcolorbox}
\tcbuselibrary{skins, breakable, theorems}

% Цвета
\definecolor{edugreen}{RGB}{85, 139, 87}    % Тёплый травянисто-оливковый
\definecolor{edublue}{RGB}{74, 119, 162}    % Чистый пастельно-голубой
\definecolor{edugray}{RGB}{125, 125, 120}   % Мягкий тёплый серый
\definecolor{eduorange}{RGB}{210, 115, 65}  % Приглушенный терракотово-оранжевый
\definecolor{edupurple}{RGB}{135, 105, 150} % Благородный черничный/лавандовый
\definecolor{edured}{RGB}{185, 80, 85}      % Мягкий бруснично-красный
\definecolor{eduteal}{RGB}{65, 135, 130}    % Спокойный морской волны/чирок
\definecolor{eduyellow}{RGB}{195, 155, 60}  % Горчично-золотой

% Макрос для создания математического блока
% Аргументы: #1=тех.имя(en), #2=отображение(ru), #3=счетчик, #4=цвет, #5=фон%, #6=шрифт, #7=префикс
\newcommand{\newmathblock}[7]{%
    \newtcbtheorem[#3]{#1}{#2}{
        enhanced, breakable, sharp corners, boxrule=0.5pt, leftrule=0pt,
        top=6pt, bottom=6pt, left=10pt, right=10pt,
        before=\vspace{0.4em}, after=\vspace{0.0em},
        attach title to upper=\par, after title={\noindent\hspace{\parindent}},
        fonttitle=\bfseries, separator sign={}, coltitle=black,
        description font=\normalfont\itshape,
        colback=#4!#5!white, colframe=#4!25, borderline west={3pt}{0pt}{#4!85},
        fontupper=#6
    }{#7}%
}

% Окружения математических блоков
\newmathblock{theorem}{Теорема}{number within=chapter}{edugreen}{4}{\normalfont}{th}
\newmathblock{lemma}{Лемма}{use counter from=theorem}{edugreen}{4}{\normalfont}{lem}
\newmathblock{definition}{Определение}{number within=chapter}{edublue}{4}{\normalfont}{def}
\newmathblock{example}{Пример}{number within=chapter}{edugray}{5}{\normalfont}{ex}

% 7. ЛОКАЛИЗАЦИЯ: Переименование стандартных названий
\addto\captionsrussian{\renewcommand{\chaptername}{Лекция}}      % Глава
\addto\captionsrussian{\renewcommand{\contentsname}{Содержание}} % Оглавление

% 8. НУМЕРАЦИЯ: Формат разделов
\renewcommand{\thesection}{§\,\arabic{section}} % Параграфы в виде: § 1

% 9. МАКРОСЫ: Сокращения для быстрого набора
\newcommand{\N}{\mathbb{N}}     % Натуральные числа
\newcommand{\Z}{\mathbb{Z}}     % Целые числа
\newcommand{\Q}{\mathbb{Q}}     % Рациональные числа
\newcommand{\R}{\mathbb{R}}     % Вещественные числа
\newcommand{\eps}{\varepsilon}  % Эпсилон

% Символы комбинаторики
\newcommand{\Ank}[2]{\mathrm{A}_{#1}^{#2}} % Размещение из n по k
\newcommand{\Cnk}[2]{\mathrm{C}_{#1}^{#2}} % Сочетание из n по k
\newcommand{\Pn}[1]{\mathrm{P}_{#1}}       % Перестановка из n

% 10. КОНЕЦ: Интерактивность и ссылки (СТРОГО В КОНЦЕ!)
\usepackage[
    colorlinks=true,   % Цветной текст для ссылок вместо цветных рамок
    linkcolor=black,   % Черный цвет ссылок в оглавлении
    urlcolor=blue,     % Синий цвет для интернет-ссылок
    citecolor=black    % Черный цвет для библиографии
]{hyperref}
